---
title: Pi4 Recovery Handler — Implementation Specification
category: algorithms
slug: pi4-recovery-handler
nav_order: 8
status: proposed
---

\section{$\Pi_4$ Recovery Handler — Implementation Specification}

\subsection*{Purpose}

This document specifies the concrete implementation required to close $\Pi_4$
in \texttt{process\_conversation} (\texttt{orchestrator\_core.py}).

$\Pi_4$ is the only hard gap between the system's current liveness guarantee
($\Diamond\ Terminated$) and full liveness ($\Diamond\ \text{Success}$).

\subsection*{The Formal Condition}

\[ \Pi_4 \triangleq \square \left( d = D_S \wedge f(C) = incomplete \rightarrow \Diamond\ RecoveryInvoked \right) \]

In implementation terms:

\begin{itemize}
  \item $d = D_S$ maps to \texttt{turn\_count >= max\_turns}
  \item $f(C) = incomplete$ maps to scratchpad containing unresolved entries
        (i.e. the supervisor did not reach a clean text completion)
  \item $RecoveryInvoked$ maps to one of three defined recovery paths
        being entered before \texttt{process\_conversation} exits
\end{itemize}

\subsection*{The Hook — Where to Intervene}

The current max-turns failsafe in \texttt{process\_conversation}:

\begin{verbatim}
# CURRENT — no recovery
if turn_count >= max_turns:
    LOG.error("ORCHESTRATOR ▸ Max turns reached.")
    await asyncio.to_thread(
        self.project_david_client.runs.update_run_fields,
        run_id,
        failed_at=int(time.time()),
        incomplete_details=(
            f"Max turns ({max_turns}) reached without clean completion."
        ),
    )
    yield json.dumps({
        "type": "error",
        "content": "Maximum conversation turns exceeded"
    })
\end{verbatim}

Replace with:

\begin{verbatim}
# PROPOSED — Pi4 recovery gate
if turn_count >= max_turns:
    LOG.error("ORCHESTRATOR ▸ Max turns reached. Invoking Pi4 recovery.")
    async for chunk in self._handle_max_turns_recovery(
        thread_id=thread_id,
        run_id=run_id,
        assistant_id=assistant_id,
        model=model,
        api_key=api_key,
    ):
        yield chunk
    return
\end{verbatim}

\subsection*{Recovery Router}

\begin{verbatim}
async def _handle_max_turns_recovery(
    self,
    thread_id: str,
    run_id: str,
    assistant_id: str,
    model: Any,
    api_key: str,
) -> AsyncGenerator[str, None]:
    """
    Pi4 Recovery Handler.

    Reads scratchpad state and routes to one of three recovery paths:
      A. PARTIAL SYNTHESIS  — verified entries exist, synthesize what is known
      B. SCOPE REDUCTION    — some verified, some pending; re-plan with reduced scope
      C. ESCALATION         — nothing verified; emit structured failure event

    Guarantees: never exits silently. Always produces either a partial
    result or an explicit structured failure signal.
    """
    scratchpad_state = await self._read_scratchpad_state(thread_id)

    verified  = scratchpad_state.get("verified",  [])
    pending   = scratchpad_state.get("pending",   [])
    failed    = scratchpad_state.get("failed",    [])

    if verified and not pending:
        # PATH A: all claims resolved — synthesize verified subset
        async for chunk in self._recover_partial_synthesis(
            thread_id, run_id, assistant_id, model, api_key,
            verified=verified, failed=failed,
        ):
            yield chunk

    elif verified and pending:
        # PATH B: partial progress — reduce scope and attempt one more delegation
        async for chunk in self._recover_scope_reduction(
            thread_id, run_id, assistant_id, model, api_key,
            verified=verified, pending=pending,
        ):
            yield chunk

    else:
        # PATH C: nothing verified — escalate, do not fabricate
        async for chunk in self._recover_escalation(
            run_id, failed=failed
        ):
            yield chunk
\end{verbatim}

\subsection*{Path A — Partial Synthesis}

\textbf{Trigger:} verified scratchpad entries exist, no pending entries remain.

\textbf{Behaviour:} construct a synthesis prompt from verified entries only.
Inject the constraint that unverified claims must be explicitly flagged.
Run one final inference turn without tool access.

\begin{verbatim}
async def _recover_partial_synthesis(
    self, thread_id, run_id, assistant_id, model, api_key,
    verified, failed,
) -> AsyncGenerator[str, None]:

    unresolved_note = ""
    if failed:
        labels = ", ".join(f['entity'] for f in failed)
        unresolved_note = (
            f"\n\nThe following could not be verified: {labels}. "
            f"Explicitly tell the user these are unresolved."
        )

    synthesis_prompt = (
        f"You have reached the maximum delegation depth. "
        f"Synthesize a final response using ONLY the verified entries "
        f"in the scratchpad. Do not infer or fabricate any values. "
        f"Cite each verified source URL."
        f"{unresolved_note}"
    )

    await self._inject_recovery_message(
        thread_id, assistant_id, synthesis_prompt
    )

    # One final stream turn — no tools, no loop
    async for chunk in self.stream(
        thread_id=thread_id,
        message_id=None,
        run_id=run_id,
        assistant_id=assistant_id,
        model=model,
        api_key=api_key,
        tool_choice="none",       # disable tool calls for this turn
    ):
        yield chunk

    await asyncio.to_thread(
        self.project_david_client.runs.update_run_fields,
        run_id,
        completed_at=int(time.time()),
        incomplete_details="Partial synthesis: max turns reached, verified subset synthesized.",
    )
\end{verbatim}

\subsection*{Path B — Scope Reduction}

\textbf{Trigger:} verified entries exist but pending entries remain unresolved.

\textbf{Behaviour:} construct a reduced-scope delegation prompt covering only
the unresolved pending items. Grant one additional turn budget ($\Delta D$).
This is the one case where the system extends beyond \texttt{max\_turns}.

\begin{verbatim}
async def _recover_scope_reduction(
    self, thread_id, run_id, assistant_id, model, api_key,
    verified, pending,
) -> AsyncGenerator[str, None]:

    pending_labels = ", ".join(p['entity'] for p in pending)
    reduction_prompt = (
        f"Maximum delegation depth reached. "
        f"The following items remain unresolved: {pending_labels}. "
        f"Attempt ONE final targeted delegation for each. "
        f"If still unresolved after one attempt, flag as [UNVERIFIED] "
        f"and proceed to synthesis. Do not loop further."
    )

    await self._inject_recovery_message(
        thread_id, assistant_id, reduction_prompt
    )

    # One additional loop iteration only
    async for chunk in self.stream(
        thread_id=thread_id,
        message_id=None,
        run_id=run_id,
        assistant_id=assistant_id,
        model=model,
        api_key=api_key,
    ):
        yield chunk

    await asyncio.to_thread(
        self.project_david_client.runs.update_run_fields,
        run_id,
        incomplete_details=(
            f"Scope reduction recovery: {len(pending)} pending items "
            f"at max turns. One additional delegation attempted."
        ),
    )
\end{verbatim}

\subsection*{Path C — Escalation}

\textbf{Trigger:} no verified entries exist. Nothing can be synthesized safely.

\textbf{Behaviour:} emit a structured escalation event. Do not fabricate.
Stamp \texttt{failed\_at}. The caller is responsible for handling escalation.

\begin{verbatim}
async def _recover_escalation(
    self, run_id, failed,
) -> AsyncGenerator[str, None]:

    failed_labels = [f.get("entity", "unknown") for f in failed]

    yield json.dumps({
        "type": "escalation",
        "reason": "max_turns_no_verified_data",
        "run_id": run_id,
        "failed_entities": failed_labels,
        "message": (
            "Research exhausted maximum delegation depth with no verified "
            "results. Manual review or query reformulation required."
        ),
    })

    await asyncio.to_thread(
        self.project_david_client.runs.update_run_fields,
        run_id,
        failed_at=int(time.time()),
        incomplete_details=(
            f"Escalation: {len(failed_labels)} entities unresolvable "
            f"after max turns."
        ),
    )
\end{verbatim}

\subsection*{Supporting Helper — Scratchpad State Reader}

\begin{verbatim}
async def _read_scratchpad_state(self, thread_id: str) -> dict:
    """
    Parse scratchpad content into structured state buckets.

    Returns:
      {
        "verified":  [{"entity": ..., "field": ..., "url": ...}],
        "pending":   [{"entity": ..., "field": ...}],
        "failed":    [{"entity": ..., "reason": ...}],
      }
    """
    raw = await self._native_exec.get_scratchpad(thread_id)
    lines = (raw or "").splitlines()

    verified, pending, failed = [], [], []

    for line in lines:
        if line.startswith("✅"):
            parts = [p.strip() for p in line.split("|")]
            if len(parts) >= 4:
                verified.append({
                    "entity": parts[1],
                    "field":  parts[2],
                    "value":  parts[3],
                    "url":    parts[4] if len(parts) > 4 else "",
                })
        elif line.startswith("🔄"):
            parts = [p.strip() for p in line.split("|")]
            if len(parts) >= 2:
                pending.append({
                    "entity": parts[1],
                    "field":  parts[2] if len(parts) > 2 else "",
                })
        elif line.startswith("⚠️") or line.startswith("❓"):
            parts = [p.strip() for p in line.split("|")]
            failed.append({
                "entity": parts[1] if len(parts) > 1 else line,
                "reason": parts[2] if len(parts) > 2 else "",
            })

    return {"verified": verified, "pending": pending, "failed": failed}
\end{verbatim}

\subsection*{Supporting Helper — Recovery Message Injector}

\begin{verbatim}
async def _inject_recovery_message(
    self, thread_id: str, assistant_id: str, content: str
) -> None:
    """
    Injects a recovery instruction into the thread as a user message.
    This is the mechanism by which the orchestrator re-steers the LLM
    without creating a new run — the next stream() call picks it up
    as the latest context.
    """
    await self._native_exec.create_message(
        thread_id=thread_id,
        assistant_id=assistant_id,
        content=content,
        role="user",
    )
\end{verbatim}

\subsection*{Formal Closure}

With this handler in place:

\begin{itemize}
  \item Path A closes $\Pi_4$ — verified subset is synthesized, unresolved claims
        are explicitly flagged. $\Diamond\ \text{Success}$ holds for the verified subset.
  \item Path B closes $\Pi_4$ — one additional targeted delegation is attempted
        before synthesis. $\Diamond\ \text{Success}$ holds with $\leq 1$ turn extension.
  \item Path C closes $\Pi_4$ — a structured escalation event is emitted.
        Silent failure is excluded. $\Diamond\ \text{Success}$ is delegated to the caller.
\end{itemize}

In all three paths:

\[ d = D_S \wedge f(C) = incomplete \rightarrow \Diamond\ RecoveryInvoked \]

$\Pi_4$ holds. Combined with $\Pi_1$ (hardened), $\Pi_2$ (implemented),
and $\Pi_3$ (infrastructure), $RPI$ holds and the Liveness Theorem is verified.

\[ \therefore \Diamond\ \text{Success} \]

\qed